\apxsection{apx-intro}{凡例}

一，本書は田川春道撰・倉田東岳画『外蕃容貌圖画』（乾坤二冊，安政二年〔一八五五〕正月刊，製本売捌 金華堂須原屋佐助）の翻刻・重排本である。底本は筑波大学附属図書館所蔵本（請求記号 ヤ600-4）で，同館「筑波大学デジタルコレクション」公開の画像による。

一，本文は原本の半丁を一頁とし，列（行）の切れ目，振り仮名，割注，圏点を原本どおりに組んだ。説明文は手書き版下の楷書に近いKlee One，序は行草書に近い衡山毛筆フォント行書を用いた。両書体にない字は全字庫正楷体で補った。

一，原本の異体字・略字（{\apxsung 叓}・国・亜・号など）は通行の旧字体（事・國・亞・號）に統一し，対照と出現数を巻末の「異体字一覧」に掲げた。誤字と思われる字（「天笠」など）は改めず，校記に記した。片仮名の「子」（ネ）は「ネ」とした。序の変体仮名は通行の平仮名に改めた。

一，判読に疑問の残る箇所，原本の誤りと思われる箇所は本文を改めず，巻末の「校記」に記した。校記の見出しは冊・画像番号・左右で示し，その頁へのリンクになっている。

一，人物図・世界図は原本の画像から紙の黄ばみと照明むらを除き，手彩色と線はそのまま残した（AIによる補筆はしていない）。図上に刷られた州名・国名，世界図の傍書は画像から除き，文字として組み直した。刊記の二印は画像から切り出した。

一，底本の蔵書印「直政之印」，図書館の登録印「明治四十一年十月五日編入」は底本個体のものであり，本書には収めない。見返しの剥離など底本の損傷も再現しない。

一，画像の利用について：筑波大学附属図書館のデジタルコンテンツ（No Copyright -- Contractual Restrictions）を，同館の利用条件（調査・研究・教育・学習目的に限る，出典の明示，改変の明示）に従って使用し，上記のとおり改変した。本書の図版を利用する場合も同じ条件に従われたい。

一，翻字と組版は Claude（Anthropic の AI）の支援を受けて行った。本文は KuroNet（人文学オープンデータ共同利用センター）の自動翻刻と照合し，振り仮名と判読に疑問のある箇所はすべて原本の画像と照合したが，専門家による校閲を経ていないため，誤りの指摘を歓迎する。
